\documentclass[11pt,letterpaper]{article}
\usepackage[utf8]{inputenc}
\usepackage[spanish,es-tabla]{babel}
\usepackage[margin=2.5cm]{geometry}
\usepackage{amsmath,amssymb}
\usepackage{booktabs}
\usepackage{graphicx}
\usepackage{float}
\usepackage{fancyhdr}
\usepackage{hyperref}
\usepackage[table]{xcolor}
\usepackage{tikz}
\usepackage{transparent}
\usepackage{eso-pic}
\usepackage{caption}
\definecolor{GeoOrange}{HTML}{FF8C00}
\definecolor{GeoBlue}{HTML}{0D47A1}
\AddToShipoutPictureBG{
  \AtPageCenter{
    \makebox[0pt]{\transparent{0.06}
\includegraphics[width=12cm]{Iconos/logo_geopol.png}}
  }
}
\hypersetup{colorlinks=true, linkcolor=GeoBlue, urlcolor=GeoOrange}
\pagestyle{fancy}
\fancyhf{}
\fancyhead[L]{\footnotesize \textcolor{GeoBlue}{\textbf{GeoPol Web}} - Reporte Técnico}
\fancyhead[R]{\footnotesize Universidad Distrital F.J.C.}
\fancyfoot[C]{\thepage}
\renewcommand{\headrulewidth}{0.4pt}
\renewcommand{\footrulewidth}{0.4pt}
\begin{document}
\begin{titlepage}
\begin{tikzpicture}[remember picture,overlay]
  \fill[GeoBlue] (current page.north west) rectangle ([yshift=-4cm]current page.north east);
  \fill[GeoOrange] ([yshift=-4cm]current page.north west) rectangle ([yshift=-4.5cm]current page.north east);
  \fill[GeoBlue!5] (current page.south west) rectangle ([yshift=4cm]current page.south east);
\end{tikzpicture}
\vspace*{-2cm}
\begin{center}
  \textcolor{white}{\Huge \textbf{GEOPORTAL WEB (GeoPol)}} \\[0.5cm]
  \textcolor{white}{\large \textit{''Máxima precisión al alcance de tus manos''}} \\[2cm]
  \vspace{3cm}
  {\LARGE \textbf{INFORME TÉCNICO DE INGENIERÍA}} \\[0.5cm]
  {\Large \textbf{Informe Técnico de Altimetría\\ (Nivelación Abierta con Control)}} \\[2cm]
  \begin{flushleft}
    \Large \textbf{Autores del Procesamiento:}\\[0.2cm]
    \large $\bullet$ Kevin Stiven Cubillos Ramirez \\[0.1cm]
    \large $\bullet$ Sergio Eduardo Barbosa Torres \\[0.1cm]
    \vspace{1cm}
    \Large \textbf{Tutor - Director de Proyecto:}\\[0.2cm]
    \large $\bullet$ Ing. Edgar Ladino
  \end{flushleft}
  \vfill
  \textbf{UNIVERSIDAD DISTRITAL FRANCISCO JOSÉ DE CALDAS}\\[0.2cm]
  Facultad de Medio Ambiente y Recursos Naturales \\ Ingeniería Topográfica / Civil \\[0.2cm]
  Bogotá D.C. -- 28 de July de 2026
\end{center}
\end{titlepage}
\tableofcontents
\newpage
\section{Marco Teórico y Referencia Altimétrica}
El presente informe detalla el cálculo, ajuste y compensación de una red de apoyo altimétrico. El proceso se fundamenta en la Nivelación Geométrica o Directa, garantizando la transferencia de cotas desde un Banco de Nivel (BM) de origen hacia los puntos de interés.
Al tratarse de una Nivelación Abierta con Control, la línea inicia en un Banco de Nivel conocido y cierra sobre un Banco de Nivel distinto, permitiendo contrastar la cota calculada de llegada con la elevación teórica esperada.
\section{Cartera Altimétrica Compensada}
A continuación se relacionan las lecturas de campo (vistas atrás, intermedias y adelante) junto con las cotas instrumentales y las elevaciones ajustadas tras el prorrateo del error de cierre:
\begin{table}[H]
  \centering
  \small
  \caption{Cartera de Nivelación Procesada (Parte 1)}
  \label{tab:nivelacion_0}
  \begin{tabular}{|c|c|c|c|c|c|}
    \hline
    \rowcolor{GeoBlue}
    \textcolor{white}{\textbf{Estaca / Punto}} & \textcolor{white}{\textbf{Vista Atrás (+)}} & \textcolor{white}{\textbf{Altura Inst. (HI)}} & \textcolor{white}{\textbf{Vista Intermedia (-)}} & \textcolor{white}{\textbf{Vista Adelante (-)}} & \textcolor{white}{\textbf{Cota Calculada}} \\
    \hline
    \rowcolor{GeoBlue!5}
    BM-INICIO & 1.500 & 501.500 &  &  & 500.000 \\
    \hline
    \rowcolor{white}
    K0+000 &  &  & 1.200 &  & 500.300 \\
    \hline
    \rowcolor{GeoBlue!5}
    PC-1 & 1.620 & 502.020 &  & 1.100 & 500.400 \\
    \hline
    \rowcolor{white}
    K0+010 &  &  & 1.450 &  & 500.570 \\
    \hline
    \rowcolor{GeoBlue!5}
    BM-LLEGADA &  &  &  & 2.505 & 499.515 \\
    \hline
  \end{tabular}
\end{table}
\vspace{0.8cm}
\begin{table}[H]
  \centering
  \small
  \caption{Cartera de Nivelación Procesada (Parte 2)}
  \label{tab:nivelacion_1}
  \begin{tabular}{|c|c|c|c|}
    \hline
    \rowcolor{GeoBlue}
    \textcolor{white}{\textbf{Estaca / Punto}} & \textcolor{white}{\textbf{Vista Atrás (+)}} & \textcolor{white}{\textbf{Corrección (m)}} & \textcolor{white}{\textbf{Cota Ajustada}} \\
    \hline
    \rowcolor{GeoBlue!5}
    BM-INICIO & 1.500 & 0.0000 & 500.000 \\
    \hline
    \rowcolor{white}
    K0+000 &  & 0.0025 & 500.303 \\
    \hline
    \rowcolor{GeoBlue!5}
    PC-1 & 1.620 & 0.0025 & 500.402 \\
    \hline
    \rowcolor{white}
    K0+010 &  & 0.0050 & 500.575 \\
    \hline
    \rowcolor{GeoBlue!5}
    BM-LLEGADA &  & 0.0050 & 499.520 \\
    \hline
  \end{tabular}
\end{table}
\vspace{0.8cm}

\section{Análisis de Errores y Compensación Altimétrica}
Las métricas de validación geométrica del circuito arrojaron los siguientes resultados:
\begin{itemize}
  \item \textbf{Sumatoria Vista Atrás ($\Sigma V^+$):} 3.120 m
  \item \textbf{Sumatoria Vista Adelante ($\Sigma V^-$):} 3.605 m
  \item \textbf{Cota Final Calculada (Cruda):} 499.515 m
  \item \textbf{Cota Teórica Esperada:} 499.520 m
  \item \textbf{Error de Cierre Altimétrico:} -0.0050 m (-5.0 mm)
\end{itemize}
\section{Perfil Topográfico de Nivelación}
\begin{figure}[H]
  \centering
  \includegraphics[width=0.95\textwidth]{Reportes_PDF/Perfil_Nivelacion.png}
  \caption{Perfil altimétrico de la línea de nivelación compensada}
\end{figure}
\section{Conclusiones}
El ajuste altimétrico se ha distribuido de manera proporcional en los puntos de cambio, obteniendo un conjunto de elevaciones definitivas aptas para la densificación de cotas, control vertical y ejecución de obras civiles.
\end{document}